\section{Does Prediction Improve Intervention Selection?}
\label{sec:selection}
The bridge motivates a CPU-only feasibility test: can a selector choose one configuration-level action that restores the fixed reference across a workload without new deviations more efficiently than direct alternatives? Reference fallback is a baseline; retaining speculation needs an operational benefit, unmeasured here.

We inspect two candidates on existing exposed evidence. A stage/state-fingerprint rule would compare aligned live state and choose an intervention at the earliest differing stage. It is not implemented: required reference/$U_2$/later native state is missing, and a stage difference alone supplies no repair direction. The same prescribed continuation change accompanies a restoration, a regression and no-ops. Post-intervention scores cannot serve as free pre-decision features.

The second candidate is generic paid proxy screening followed by actual-action verification. An offline query interface reveals retained outcomes only after charging their acquisition. The policy sees handles and paid replies without task/platform names, reference tokens or unqueried labels. Accepted actions require verification; unknown counterfactuals remain unknown. Both orders of the two RTX actions and the sole B570 action are compared with direct enumeration; no per-task mixture is selected after seeing outcomes.
\begin{table}[t]
\centering
\caption{Offline selection replay after common reference/baseline acquisition. Counts are available historical record equivalents, not measured time or new inference.}
\label{tab:selection}
\begin{tabular}{@{}llrl@{}}
\toprule
Workload & Policy & Records & Decision \\
\midrule
RTX19 & Enumerate & 114 & Abstain \\
RTX19 & Proxy + verify & 114 & Abstain \\
B57024 & Verify native & 72 & Select $N_1$ \\
B57024 & Proxy + verify & 144 & Select $N_1$ \\
\bottomrule
\end{tabular}
\end{table}


On RTX, screening $N_1$ through its proposal-conditioned predictor and then testing $U_2$ leads to abstention, as do both enumeration orders. On B570, screening through $T_1$ and then verifying $N_1$ selects the same sole candidate as direct verification but adds 72 proxy records. These are retrospective acquisition equivalents, not new model executions, independent cases or latency measurements. The already-exposed MBPP supplement has two persistent failures and no observed repairs; its entire $U_2$ counterfactual remains unknown.

For a sole accepted action, enumeration costs $O+N+V$ and screening costs $O+P+A+N+V$, where $O$ is common observation, $P$ proxy work, $A$ extra selection/setup and $V$ verification outside the candidate execution. Screening adds $P+A$. For a rejected action, savings require $P+A+(V_P-V_E)<N$, at a comparable success and false-rejection contract. RTX prediction retains substantial numerical work; no matched timing establishes that inequality. Cached decisions can also benefit enumeration. No debugging-time or deployed-cost saving follows.

The result is a scoped negative feasibility finding: neither candidate supplies a supported distinct selection capability here. Ten integrity/charging tests pass, without held-out validation. The supplement retains ledgers, missingness and cost cases.
